\chapter*{General Introduction}
\addcontentsline{toc}{chapter}{General Introduction}
\thispagestyle{fancy}
\justifying

The effectiveness of emergency medical systems is a critical determinant of public health outcomes, particularly in developing nations like Burkina Faso. At the heart of these systems lies the SAMU (Service d'Aide M\'edicale Urgente), an institution dedicated to providing life-saving pre-hospital care. However, the success of any emergency response is profoundly dependent on the initial sorting and prioritization of patients, a process formally known as triage. Historically, this essential task has been performed manually within the Burkinab\`e context, relying on verbal communication and the subjective experience of healthcare providers. While this traditional approach has served its purpose, it faces significant limitations in an era of increasing emergency volumes and complex clinical scenarios.

\section*{Problem Statement}
\addcontentsline{toc}{section}{Problem Statement}
In the current operational framework of the SAMU Burkina Faso, triage decisions are frequently made under extreme pressure, often based on incomplete or imprecise information provided by callers in a state of distress. This reliance on human intuition without standardized digital assistance introduces several vulnerabilities. Geographical constraints, such as urban traffic and poor road conditions in rural areas, already impose inherent delays on interventions. When these are compounded by potential inconsistencies in patient prioritization, the risk to life increases significantly. The absence of a structured, digital decision support tool means that the most critical cases may not always be identified with the necessary speed and precision, leading to suboptimal allocation of the SAMU's specialized mobile medical units (SMUR). Consequently, there is an urgent need to modernize this process through a technological solution capable of formalizing and securing triage decisions.

\section*{Objectives}
\addcontentsline{toc}{section}{Objectives}
The primary goal of this research is to design and develop a comprehensive digital decision support system for pre-hospital triage at the SAMU of Burkina Faso. This system is fundamentally based on the internationally validated Emergency Severity Index (ESI) algorithm, which provides a clinically rigorous framework for patient sorting. To achieve this, the project focuses on several specific objectives. First, we aim to analyze and define the functional requirements specific to the Burkinab\`e pre-hospital environment. Second, the project involves implementing an automated decision tree that digitizes the 5-level ESI protocol. Third, we seek to develop a secure, high-performance full-stack web application that integrates these functionalities for both frontline operators and medical supervisors. Finally, our objective includes providing advanced dashboarding and reporting features to ensure real-time visibility and auditability of triage activities.

\section*{Methodology}
\addcontentsline{toc}{section}{Methodology}
To ensure the successful realization of this project, we have adopted the Two-Track Unified Process (2TUP) as our primary development methodology. This approach is particularly well-suited for complex software projects as it allows for the parallel management of functional and technical requirements. The functional track of our work focuses on the clinical logic of the ESI algorithm and the specific workflows of the SAMU agents. Simultaneously, the technical track addresses the implementation of a modern architecture, utilizing FastAPI for the backend and React.js for the frontend, while ensuring robust data security through PostgreSQL and advanced authentication mechanisms. The convergence of these two tracks results in a solution that is both clinically accurate and technically resilient.

\section*{Thesis Organization}
\addcontentsline{toc}{section}{Thesis Organization}
This thesis is structured into two main parts, each addressing a critical dimension of the project. \textbf{Part I} establishes the theoretical foundations and reviews the state of the art in the field of digital health and emergency triage. Chapter 1 explores the essential concepts of emergency medical systems and the ESI model, while Chapter 2 details the methodology and the conceptual approach adopted. \textbf{Part II} focuses on the design and implementation of the proposed solution. Chapter 3 presents the technical modelling and the development process, and Chapter 4 analyzes the results, functional performance, and future perspectives for the system. A general conclusion synthesizes our contributions and outlines further improvements.

\clearpage
